\documentclass{article}
\usepackage[utf8]{inputenc}
\usepackage{amsmath}

% Dies ist ein Template für KantiKoala Übungen.
% Verwende die 'question' Umgebung für Aufgaben.

% Format:
% \begin{question}{DATENTYP}{LÖSUNG}
%    Fragetext...
%    \item A) Option 1 (nur bei Multiple Choice)
%    \item B) Option 2
%    \explanation{Erklärung für die Lösung...}
% \end{question}

\newenvironment{question}[2]{
    \textbf{Frage (#1):} 
    % Im PDF wird dies als normaler Text angezeigt.
    % KantiKoala wird diesen Block extrahieren und interaktiv machen.
}{}

\newcommand{\explanation}[1]{\textit{#1}}

\begin{document}

\section*{Übungen zu Thema X}

% Beispiel 1: Multiple Choice
\begin{question}{multiple-choice}{B}
    Was ist die Hauptstadt von Frankreich?
    \item A) London
    \item B) Paris
    \item C) Berlin
    \explanation{Paris ist seit Jahrhunderten das politische Zentrum.}
\end{question}

% Beispiel 2: Textaufgabe
\begin{question}{text}{12}
    Berechne: $3 \times 4$
    \explanation{Einmaleins!}
\end{question}

\end{document}
